\newpage
\section{Drehzahlmessung mit Gabellichtschranke und Reflexoptokoppler}

\subsection{Aufgabenstellung}

Mit dem Timer NE555 ist bei einer Betriebsspannung von \(U_B = 5\,\text{V}\) ein Monoflop zur Frequenz-/Spannungsumsetzung eines Reflexoptokopplersignals (siehe Abbildung 6.11) aufzubauen. Der Fototransistor des Reflexoptokopplers ist mit einem Pull-up-Widerstand \(R_1 = 100\,\text{k}\Omega\) bei \(V_{CC}=U_B\) zu betreiben.

Der Widerstand \(R\) des Monoflops ist für \(C = 1\,\mu\text{F}\) so zu dimensionieren, dass die Schaltung für Periodendauern des Reflexoptokopplersignals im Bereich
\[
	10\,\text{ms} \leq T_{\text{Reflexion}} \leq 600\,\text{ms}
\]
funktioniert.

Zur Glättung der Monoflop-Ausgangsspannung ist ein RC-Tiefpass vorzusehen. Dessen Zeitkonstante \(\tau = R_{\text{TP}}C_{\text{TP}}\) ist so zu wählen, dass der Spannungsrippel von \(U_A\) kleiner als \(0{,}2\,\text{V}\) bleibt (Richtwert: \(C_{\text{TP}} > 1\,\mu\text{F}\)). Die untere messbare Drehzahlgrenze ist experimentell zu bestimmen. Die Signalverläufe \(U_{\text{TH}}\), \(U_{\text{TRIG}}\), \(U_{\text{OUT}}\) und \(U_A\) sind beispielhaft zu skizzieren.

Zur Untersuchung des Übertragungsverhaltens ist die Ausgangsspannung \(U_A\) mit einer Dreiecksspannung des Funktionsgenerators aufzunehmen (\(V_{pp}=4\,\text{V}\), \(V_{\text{Offset}}=2\,\text{V}\), \(T=100\,\text{s}\)). Es ist zu prüfen, ob ein lineares \(f/U_A\)-Verhalten vorliegt, in welchem Drehzahlbereich dieses gegeben ist und wodurch der lineare Bereich begrenzt wird. Dazu sind Messungen mit unterschiedlichen Widerstandswerten \(R\) durchzuführen.

\subsection{Aufbau und Durchführung}
% TODO

\begin{figure}[h]
	\centering
	\begin{circuitikz}[transform shape, scale=0.8]

		% Gemeinsame Masse
		\node[rground] (GND) at (11,-1.5) {};

		% --- LED / Fotodiode ---
		\draw
		(1,4.5) node[vcc]{$U_B$}
		to[R=$R_v$] (1,2.5)
		-- (1,2)
		to[led] (1,0);
		\node[rground] at (1,0) {};

		\node[left] at (0.7,1.0) {$R_d$};

		% --- Fototransistor Q1 ---
		\draw
		(3,4.5) node[vcc]{$U_B$}
		to[R=$R_1$] (3,3) coordinate (NTRIG);
		\draw
		(3,2) node[npn, photo, anchor=C] (Q1) {};
		\draw (Q1.E) -- (3,0);
		\node[rground] at (3,0) {};
		\node[right] at (Q1) {$Q_1$};
		\draw (Q1.C) -- (NTRIG);

		% --- NE555 Block ---
		\draw (9,0) rectangle (13,5);
		\node at (11,2.5) {$NE555$};

		% Pinnummern am Gehäuse
		\node[anchor=north] at (10,5) {4};
		\node[anchor=west] at (9,3) {2};
		\node[anchor=west] at (9,2) {3};
		\node[anchor=south] at (11,0) {1};
		\node[anchor=north] at (12,5) {8};
		\node[anchor=east] at (13,3) {7};
		\node[anchor=east] at (13,2) {6};
		\node[anchor=east] at (13,1) {5};

		% Trigger (Pin 2) von Q1
		\draw (NTRIG) -- (9,3);
		\node[above] at (6.5,3) {Trigger};

		% Reset (Pin 4)
		\draw (10,6) node[above]{Reset} -- (10,5);
		\node[ocirc] at (10,6) {};

		% Output-Pin (3)
		\coordinate (OUTPIN) at (9,2);

		% Versorgung Pin 8 + Widerstand R
		\coordinate (VCCPIN) at (12,5);
		\draw (VCCPIN) -- (12,5.8);
		\node[vcc] (VB) at (12,5.8) {$U_B$};
		\draw (VB) -- (15,5.8);
		\draw (15,5.8) to[R=$R$] (15,3.5);

		% Bus für Discharge (7) und Threshold (6)
		\draw (15,3.5) -- (15,0.2);
		\draw (13,3) -- (15,3);   % Pin 7
		\draw (13,2) -- (15,2);   % Pin 6
		\node[right] at (15,3) {Discharge};
		\node[right] at (15,2) {Threshold};

		% Timing-Kondensator C zum Massepunkt (schmaler, Text links)
		\draw (15,1)
		to[C, l_=$C$] (15,-1.0) -- (15,-1.5) -- (GND);
		\node[right, align=left] at (15.6,0.0) {Control\\voltage};

		% Pin 5 mit 10 nF nach Masse (schmaler)
		\draw (13,1) -- (13.6,1) -- (13.6,-0.5)
		to[C, l_=$10\,\mathrm{nF}$] (13.6,-1.0) -- (13.6,-1.5) -- (GND);

		% Masse Pin 1
		\draw (11,0) -- (GND);

		% --- Ausgangsnetzwerk U_A, R_TP, C_TP ---
		% Knoten U_A (elektrischer Knoten bleibt bei x=6)
		\coordinate (UA) at (6,2);

		% C_TP von U_A nach Masse (bleibt an derselben Stelle)
		\coordinate (UA_C_bot) at (6,-0.8);
		\draw (UA) to[C=$C_{TP}$] (UA_C_bot);
		\node[rground] at (UA_C_bot) {};

		% R_TP zwischen U_A und Output-Pin
		\draw (UA) to[R=$R_{TP}$] ($(OUTPIN)+(-1,0)$) -- (OUTPIN);
		\node[above] at (8.3,2) {Output};

		% U_A-Spannungspfeil links von C_TP mit Abstand 1
		\coordinate (UAmeasTop) at ($(UA)+(-1,0)$);       % (5,2)
		\coordinate (UAmeasBot) at ($(UA_C_bot)+(-1,0)$); % (5,-0.8)

		\node[ocirc] (UAnode) at (UAmeasTop) {};
		\node[ocirc] (UAref)  at (UAmeasBot) {};
		\draw (UAnode) -- (UA);

		% Spannungs-Pfeil U_A mit eigenem rground
		\draw[->] ($(UAnode)+(0,-0.2)$) -- ($(UAref)+(0,0.2)$) node[midway,left] {$U_A$};
		\node[rground] at (UAref) {};

	\end{circuitikz}
	\caption{Schaltung mit Gabellichtschranke und NE555}
\end{figure}

\subsection{Dimensionierung}
% TODO 


\subsection{Durchführung}
% TODO

\subsection{Ergebnisse}
% TODO

\subsection{Erkenntnisse}

% TODO
